\documentclass[11pt]{article}
\usepackage[a4paper,margin=0.9in]{geometry}
\usepackage{amsmath,amssymb,amsthm}
\usepackage{graphicx}
\usepackage{booktabs,longtable,array,tabularx}
\usepackage{caption}
\usepackage{enumitem}
\usepackage{microtype}
\usepackage{xcolor}
\usepackage[colorlinks=true,allcolors=blue!45!black]{hyperref}
\usepackage{pdflscape}
\setlength{\tabcolsep}{5pt}
\setlength{\LTpre}{6pt}
\setlength{\LTpost}{6pt}
\setlist{nosep,leftmargin=1.5em}
\emergencystretch=3em
\graphicspath{{figs/}}
\newtheorem{proposition}{Proposition}
\newtheorem{lemma}{Lemma}
\newtheorem{remark}{Remark}
\newcommand{\Tr}{T_{\mathrm r}}
\newcommand{\Rho}{\rho}
\newcommand{\Emax}{E_{\max}}
\begin{document}
\title{Sample-and-hold governance: review intervals as a design variable in resource stability}
\author{Amin Abaee\\[0.35em]
{\small Independent Researcher}\\[0.55em]
{\small\href{https://orcid.org/0000-0002-0019-1842}{ORCID: 0000-0002-0019-1842}}\\[0.3em]
{\small\href{mailto:amin\_abaee@ut.ac.ir}{amin\_abaee@ut.ac.ir}}}
\date{October 16, 2026}
\maketitle

\begin{abstract}
Institutions observe renewable resources at discrete times, revise rules at scheduled reviews, and hold commands between them. We model this architecture as sample-and-hold governance and distinguish it from both continuous-delay feedback and an annual-step approximation. We prove forward invariance of the nonnegative state space and finite-horizon consistency as the review interval tends to zero; neither result implies stability at any fixed positive interval. On an illustrative logistic core, the exponential frozen-assessment update has one unit-circle crossing at 6.501 years, whereas a forward-Euler command update reports crossings at 47.536 and 79.143 years. The exponential protective comparator is stable on the tested interval grid. The 6.501-year value is not a calibrated institutional threshold: across six-parameter joint perturbation boxes, a crossing is found at 60 of 64 corners at \(\pm0.5\%\) and 44 of 64 at \(\pm1\%\); every no-crossing corner is stable on a 20,001-point scan. In a 42-stock screen, the Benjamini--Hochberg-adjusted target-band count is zero. A recovered review-interval proxy test has only two proxy values, both below the modeled response windows, and its negative association is opposite the predicted approach gradient and confounded by region. A 32-system case inventory identifies no candidate meeting all four eligibility criteria. Northern cod contributes a descriptive two-window partition, not an identified governance mechanism. The results support operator-specific stability analysis, not a general claim that a particular review cadence stabilizes real fisheries.
\end{abstract}

\noindent\textbf{Keywords:} sample-and-hold control; review interval; resource governance; sampled-data stability; fisheries; management strategy evaluation.

\noindent\textbf{Evidence status.} The analytic propositions, deterministic multiplier calculations, archived trajectory classifications, empirical screens, and prospective designs have distinct evidential status. Supplementary Table S1 inventories the claims and limitations; parameter values in the numerical core are illustrative and uncalibrated.

